\documentclass[11pt,letterpaper]{article}
\usepackage[T1]{fontenc}
\usepackage[utf8]{inputenc}
\usepackage{lmodern,microtype}
\usepackage[margin=1in]{geometry}
\usepackage{amsmath,amssymb,amsthm,mathtools,mathrsfs}
\usepackage{booktabs,tabularx,array,longtable,enumitem}
\usepackage{xurl}
\usepackage[hidelinks]{hyperref}
\usepackage{tocloft}
\hypersetup{pdftitle={Fermionic Theta Transfers, Open Quantum Dynamics, and Energy Response},pdfauthor={Drafted for Edward Baker}}
\setlength{\parskip}{3pt}
\setlength{\emergencystretch}{2em}
\setlength{\cftbeforesecskip}{2pt}
\setcounter{tocdepth}{1}
\setlist{itemsep=3pt,topsep=5pt}
\numberwithin{equation}{section}
\newtheorem{theorem}{Theorem}[section]
\newtheorem{proposition}[theorem]{Proposition}
\newtheorem{lemma}[theorem]{Lemma}
\theoremstyle{remark}
\newtheorem{remark}[theorem]{Remark}
\DeclareMathOperator{\Tr}{Tr}
\DeclareMathOperator{\Var}{Var}
\DeclareMathOperator{\Cov}{Cov}
\DeclareMathOperator{\diag}{diag}
\DeclareMathOperator{\supp}{supp}
\newcommand{\E}{\mathbb E}
\newcommand{\N}{\mathbb N}
\newcommand{\dd}{\,d}
\newcommand{\ket}[1]{\lvert #1\rangle}
\newcommand{\bra}[1]{\langle #1\rvert}
\newcommand{\ip}[2]{\langle #1,#2\rangle}
\newcommand{\loss}{\mathsf R}
\newcommand{\srcR}{\mathscr R}
\newcommand{\srcS}{\mathscr S}
\newcommand{\ev}[1]{\textsf{[#1]}}
\newcommand{\Fock}{\mathcal F_-}
\title{Fermionic Theta Transfers, Open Quantum Dynamics,\\and Energy Response\\[6pt]
\large A Continuation of the Signed Mellin-Curvature Program}
\author{Drafted for Edward Baker}
\date{2 November 2026}
\begin{document}
\maketitle
\begin{abstract}
We continue the background manuscript on signed Mellin curvature and theta
coefficient positivity. The coefficient determinants are encoded as signed
fermionic amplitudes of a normalized backward-shift transfer. We summarize
the complete-source estimates proving that its activity generator is
dissipative on $0\le v\le1/25$, construct the associated particle-loss
count instrument, and state the fixed-mode Hermitian collision error and
the compatible vacuum-Fock extension. A coherent reference recovers the
amplitude, including its sign, with explicit success normalizations.
These realization results do not provide the required determinant
comparison. A worked obstruction shows why preservation of the fixed
ordered cone on every intermediate interval would force a Gaussian
generating function, while the target ratio is exactly independent of
activity. We then connect the derivative orbitals to the natural theta
energies $E_m=\pi m^2$: each orbital samples a filtered thermal response,
rather than representing one energy eigenstate. The remaining analytical
obligation is a source-dependent bound on the joint signed response. No
new range of the Gaussian Mellin comparison or of coefficient positivity
is established by this synthesis.
\end{abstract}

\noindent\textbf{Reading contract and status.}
This paper assumes the setup and signed-averaging baseline of
\emph{Signed Mellin Curvature and Theta Coefficient Positivity:
Mathematical Background and Physical Research Directions} (\ev{BG}).
It develops the fermionic, open-system and energy interfaces in detail,
rather than repeating the background's survey of all eight branches.
The local realization has a separate same-model internal LLM assessment;
the thermal and original fermionic dictionaries have cross-model internal
assessments. The intermediate-cone obstruction has proof-author checking
only, with its separate check still pending. None is independent
specialist validation. Appendix~\ref{sec:evidence} identifies the source
records and their precise status.

\clearpage
\tableofcontents
\clearpage

\section{The unchanged mathematical target}
\label{sec:setup}
We retain the background's complete source at physical heat time zero:
\begin{align}
 p(u)&=e^u\sum_{m\ge1}e^{-\pi m^2e^{4u}},\qquad
 \Phi(u)=\frac{p''(u)-p(u)}8,\notag\\
 c_n&=\frac1{(2n)!}\int_0^\infty u^{2n}\Phi(u)\dd u,
 \qquad F(w)=\sum_{n\ge0}c_nw^n. \label{eq:source}
\end{align}
Here $c_n=0$ for negative integers. The source is positive on the
half-line, even after full theta reconstruction, and has
 double-exponential decay. In particular $c_n>0$ for $n\ge0$, $F$ is
entire, and the background normalization is
$F(z^2)=\xi((1+z)/2)/8$. Write
\begin{equation}
 D_{r,k}=\det[c_{k+j-i}]_{i,j=0}^{r-1},\qquad
 M_3(k)=(2k)!(2k+2)!(2k+4)!D_{3,k}. \label{eq:det}
\end{equation}
The background supplies $D_{3,k}>0$ for every integer $k\ge0$, and
positivity of its signed real-Mellin interpolation $M_3(s)$ for $s\ge0$.
These are inherited premises, not consequences of a positive quantum
channel.

The approach separates three questions: can the coefficient transfer be
realized by physical evolution, can its signed amplitude be recovered by
an observation, and does the source impose a useful inequality on that
amplitude? A quantum channel acts on density matrices and must be
completely positive and trace preserving (CPTP). A count instrument also
records detected losses. Conditioning on no count retains a linear
amplitude operator that can be non-Hermitian; its entries can have either
sign. The construction below answers the realization and readout
questions while leaving the source-dependent comparison open.

The finite-order quantitative consumer is
\begin{equation}
 \eta_k:=\frac{D_{3,k-1}D_{3,k+1}}{D_{3,k}^2}
 \ \le\ \frac{k}{k+3},\qquad k\ge1. \label{eq:target}
\end{equation}
Desnanot--Jacobi then gives
\begin{equation}
 D_{4,k}D_{2,k}=D_{3,k}^2(1-\eta_k),\qquad
 D_{4,k}\ge\frac{3}{k+3}\frac{D_{3,k}^2}{D_{2,k}}>0.
 \label{eq:payoff}
\end{equation}
The background explains that low-order positivity is already known;
this is a method pilot seeking a quantitative mechanism that might
extend to higher orders. Nonnegativity of \emph{all} finite Toeplitz
minors is the inherited RH-equivalent destination, not a result proved
by the present physical realization.

A sufficient route to \eqref{eq:target} is the background conjecture
\begin{equation}
 (\log M_3)''(s)\le\kappa_3(s),\qquad
 \kappa_3(s)=\sum_{j=0}^2\sum_{\nu=0}^\infty
 (s+j+\tfrac12+\nu)^{-2},\qquad s\ge0. \label{eq:curvtarget}
\end{equation}
Triangular integration over $[k-1,k+1]$ and the factorials in
\eqref{eq:det} give \eqref{eq:target}. The pointwise bound is stronger
than the required integrated comparison; neither is obtained below.
In particular, the first neighboring-index problem uses $[0,2]$.

\begin{table}[htbp]
\centering\small
\renewcommand{\arraystretch}{1.16}
\begin{tabularx}{\textwidth}{@{}>{\raggedright\arraybackslash}p{.18\textwidth}X@{}}
\toprule
Symbol & Meaning and distinction\\
\midrule
$u$, $\beta=e^{4u}$ & Source coordinate and inverse temperature.\\
$s$ & Real Mellin tilt; the source couples to $2\sum_j\log u_j$.\\
$v$ & Transfer activity; also the chosen clock for the constructed loss dynamics.\\
$i,j$ & Derivative/moment labels, or coefficient sites in the discrete model.\\
$m$, $E_m$ & Theta spectral mode and its energy $\pi m^2$.\\
$r$, $N$ & Particle number and largest retained coefficient-site index.\\
$\srcR$, $\srcS$ & Scalar ratios $\Phi''/\Phi$ and $\Phi^{(4)}/\Phi$; $\loss$ below is the loss operator.\\
$\lambda_n$ & Coefficients of $\log(F/c_0)$, denoted $b_n$ in the MC08 notes; they are not the background's raw moments $b_n$.\\
\bottomrule
\end{tabularx}
\caption{The parameters and indices are not identified by the physical interpretation.}
\label{tab:notation}
\end{table}

\section{From theta coefficients to fermionic amplitudes}
\label{sec:fermions}

\subsection{Continuum orbitals and the role of their labels}
On $L^2((0,\infty),du)$ define
\begin{equation}
 f_i=\frac{\Phi^{(2i)}}{\sqrt\Phi},\qquad
 g_j=\frac{u^{2j}\sqrt\Phi}{(2j)!},\qquad i,j\ge0.
 \label{eq:orbitals}
\end{equation}
The complete-source derivative bounds, together with its positive first
summand, make these fixed-order orbitals square integrable. Integration by parts, including the vanishing of
all odd source derivatives at zero, gives
\begin{equation}
 \ip{f_i}{g_j}=c_{j-i}. \label{eq:overlap}
\end{equation}
For $j<i$, transferring the polynomial derivatives first leaves an odd
source derivative at the origin, hence zero. This retains the small-index
boundary convention rather than extending a factorial formula beyond
its domain.

Let $a(f)$ and $a^\dagger(f)$ satisfy the canonical anticommutation
relations (CAR), and fix increasing creation order. The determinant is
\begin{equation}
 D_{r,k}=\bra{\Omega}a(f_{r-1})\cdots a(f_0)
 a^\dagger(g_k)\cdots a^\dagger(g_{k+r-1})\ket{\Omega}.
 \label{eq:car}
\end{equation}
Equivalently it is the overlap of two different exterior products.
It is not a squared norm. Positivity of the Hilbert-space inner product
controls its magnitude by Cauchy--Schwarz, but does not give its sign.
The CAR identity is standard; the complete theta orbitals and endpoint
payments are the local construction \ev{F1,F2}.

No Hamiltonian was assigned to the labels $i,j$ in
\eqref{eq:orbitals}. They record which derivative and moment occur.
The families are not orthonormal energy bases. Assigning distinct
eigenvalues to nonorthogonal orbitals would not define a self-adjoint
Hamiltonian. Section~\ref{sec:energy} gives their actual energy-response
interpretation.

\subsection{One discrete transfer for every finite minor}
Let $\mathfrak h=\ell^2(\N_0)$, $S\ket0=0$, and
$S\ket j=\ket{j-1}$ for $j\ge1$. For $v\ge0$ set
\begin{equation}
 K_v=\frac{F(vS)}{F(v)},\qquad K_0=I. \label{eq:Kv}
\end{equation}
Since $\|S\|=1$ and $c_n\ge0$,
$\|F(vS)\|\le F(v)$; thus $K_v$ is a contraction for every $v\ge0$.
For increasing finite occupation lists $I=(i_1,\ldots,i_r)$ and
$J=(j_1,\ldots,j_r)$, exterior-power expansion yields
\begin{equation}
 A_v(I,J):=\bra I\Gamma(K_v)\ket J
 =\frac{v^{\sum_bj_b-\sum_ai_a}}{F(v)^r}
       \det[c_{j_b-i_a}],\qquad v>0. \label{eq:minoramp}
\end{equation}
Here $\Gamma(K)=\bigoplus_{r\ge0}\wedge^rK$ on vacuum Fock space,
with vacuum block $+1$. Formula \eqref{eq:minoramp} has a positive
prefactor. The cases $v=0$ are read directly from $K_0=I$, without
ambiguous powers of zero. The continuum and discrete constructions
encode the same pairings; this is not an asserted unitary change of
basis between their orbital families.

Already on $\operatorname{span}\{\ket0,\ket1\}$,
\begin{equation}
 K_v=\frac1{F(v)}
 \begin{pmatrix}c_0&vc_1\\0&c_0\end{pmatrix},\qquad v>0. \label{eq:jordan}
\end{equation}
The nonzero Jordan block cannot be self-adjoint in any positive definite
inner product: similarity by the square root of such a metric would
make it Hermitian and therefore diagonalizable. A global positive metric
making $K_v$ self-adjoint would restrict to one on this invariant block.
This excludes that interpretation of $K_v$; it does not exclude a
Hermitian system--environment realization. Endpoint contractivity alone
also does not prove that the activity evolution is dissipative at every
intermediate time.

\section{An explicit interval for the exact activity generator}
\label{sec:local}

\subsection{Bounds from the complete source and an analytic logarithm}
The explicit interval comes from elementary bounds that retain the
entire theta sum and integration tail (\ev{O1,O3}). Put
\begin{equation}
 J_0=\int_0^\infty p(u)\dd u,\quad
 J_2=\int_0^\infty u^2p(u)\dd u,\quad
 C=\int_0^\infty p(u)\cosh(2u)\dd u.
\end{equation}
Full theta inversion gives $p'(0)=-1/2$, and integration by parts gives
\begin{equation}
 c_0=\frac1{16}-\frac{J_0}8,\qquad
 c_1=\frac{2J_0-J_2}{16},\qquad
 F(4)=\frac1{16}+\frac{3C}8. \label{eq:momconstants}
\end{equation}
The estimates used are
\begin{equation}
 \frac1{1600}<J_0<\frac1{209},\qquad
 J_2<\frac2{25289},\qquad C<\frac{11}{2223}. \label{eq:Jbounds}
\end{equation}
For the upper bounds, $e^{4u}\ge1+4u$ gives
$p(u)\le\sum_{m\ge1}e^{-\pi m^2}e^{-(4\pi m^2-1)u}$.
Use $4\pi m^2-1>11$ and
$\sum_{m\ge1}e^{-\pi m^2}<\sum_{m\ge1}20^{-m^2}\le1/19$.
The elementary Laplace integrals are $1/a$, $2/a^3$, and
$a/(a^2-4)$. For the lower bound, retain $m=1$ on $[0,1/16]$,
where $\pi e^{4u}<9/2$ and $e^{9/2}<100$.
The rational exponential bounds and all fractions were checked in the
existing E01 support record; no numerical integration was used.

It follows that
\begin{equation}
 \frac{207}{3344}<c_0<\frac1{16},\quad
 \lambda_1:=\frac{c_1}{c_0}>\frac1{1000},\quad
 q_4:=\frac{F(4)-c_0}{c_0}<\frac{320}{8073}<\frac1{25}.
 \label{eq:cauchyinput}
\end{equation}
Positive coefficients imply $|F(w)/c_0-1|\le q_4$ on $|w|\le4$.
The analytic logarithm with value zero at the origin therefore satisfies
\begin{equation}
 B(w):=\log(F(w)/c_0)=\sum_{n\ge1}\lambda_nw^n,
 \quad |B(w)|<\frac1{24},\quad
 |\lambda_n|\le\frac1{24\,4^n}. \label{eq:logbound}
\end{equation}
There is a genuine analytic neighborhood of this closed disk. The
estimate is not a global zero-free assertion about $F$.

\subsection{The loss operator and its gapless lower bound}
Norm-convergent power series give, for $0\le v\le1/25$,
\begin{align}
 G(v)&=K'_vK_v^{-1}
 =SF'(vS)F(vS)^{-1}-\gamma(v)I,\notag\\
 &=\sum_{n\ge1}\alpha_n(v)(S^n-I),\qquad
 \alpha_n(v)=n\lambda_nv^{n-1},\quad \gamma=F'/F. \label{eq:G}
\end{align}
In particular $G(0)=\lambda_1(S-I)$. Define
\begin{equation}
 \loss(v)=-(G+G^\dagger),\qquad
 H(v)=\frac i2(G-G^\dagger),\qquad G=-iH-\loss/2. \label{eq:HR}
\end{equation}

\begin{theorem}[Local admissibility, \ev{O1,O3}]
On the infinite one-particle space,
\begin{equation}
 \loss(v)\ge\frac1{8000}(2I-S-S^\dagger)\ge0,
 \qquad 0\le v\le\frac1{25}. \label{eq:lossbound}
\end{equation}
Writing $A_*=[96(99/100)^2]^{-1}$, one also has
\begin{equation}
 \|G\|,\|H\|\le2A_*,\quad \|\loss\|\le4A_*,\quad
 \|K_v^{-1}\|<13/12. \label{eq:normbounds}
\end{equation}
The operators $G,H,\loss$ and $\sqrt\loss$ are norm-continuous. Every
intermediate $V(b,a)=K_bK_a^{-1}$, $0\le a\le b\le1/25$, is a contraction.
\end{theorem}
\begin{proof}[Proof of the decisive bound]
The Toeplitz quadratic form of $\loss$ has continuous symbol
\[
 r_v(\theta)=2\sum_{n\ge1}\alpha_n(v)(1-\cos n\theta).
\]
For a finite sequence $x$, integrate this symbol against
$|\sum_jx_je^{ij\theta}|^2/(2\pi)$; norm convergence extends the identity
to $\ell^2$. The bound
$1-\cos n\theta\le n^2(1-\cos\theta)$ yields
\[
 \loss(v)\ge\left(\lambda_1-\sum_{n\ge2}n^3|\lambda_n|v^{n-1}\right)
 (2I-S-S^\dagger).
\]
With $x=v/4\le1/100$, the coefficient tail is at most
\[
 \frac1{96}\left(\frac{1+4x+x^2}{(1-x)^4}-1\right)
 \le\frac{2650133}{3073907232}<\frac7{8000}.
\]
Combine this with $\lambda_1>8/8000$. The same majorant controls the
zero-angle quotient; no angle sampling or positive spectral gap is
needed. Summing the coefficient norms gives \eqref{eq:normbounds};
the inverse follows from the Neumann series and
$(1+q_4)/(1-q_4)<13/12$. Polynomial approximation to $\sqrt{x}$ on a
common bounded spectral interval proves square-root continuity.
Finally
$\partial_b\|V(b,a)f\|^2=-\ip{V(b,a)f}{\loss(b)V(b,a)f}\le0$.
\end{proof}

The spectral infimum of $\loss$ is nevertheless zero: $r_v(0)=0$,
and normalized Hardy kernels concentrate its quadratic form near that
angle. The bound is relative to a gapless shift Laplacian. The chosen
interval is sufficient, not asserted optimal, and it does not cover
$v=1$.

There is already a bounded one-particle CPTP completion with jump
$\sqrt\loss$ acting \emph{within} $\mathfrak h$:
\[
 \dot\rho=G\rho+\rho G^\dagger+\sqrt\loss\,\rho\sqrt\loss.
\]
Its convergent time-ordered jump series proves complete positivity,
and $G+G^\dagger+\loss=0$ proves trace preservation. This is different
from the particle-loss construction below.

\subsection{What local admissibility does not imply}
The Cauchy argument also gives some sufficiently small admissible
interval for any real analytic positive-coefficient source with
$c_0,c_1>0$ and a convergent coefficient sum on a disk. The explicit
interval above is paid for theta, but the qualitative local property
is not an arithmetic discriminator.

A decisive control is $F_*(w)=1+w+2w^2$. For $x=\cos\theta$ its symbol is
\begin{equation}
 r_{*,v}(\theta)=
 \frac{2(1-x)(1-2v^2)(1+7v+2v^2+8vx)}
 {(1+v+2v^2)|1+ve^{i\theta}+2v^2e^{2i\theta}|^2}\ge0,
 \quad 0\le v<1/\sqrt2. \label{eq:controlsymbol}
\end{equation}
Yet its two-particle endpoint amplitude is
\begin{equation}
 \bra{0,1}\Gamma(K_{*,v})\ket{1,2}
 =-\frac{v^2}{(1+v+2v^2)^2}<0. \label{eq:negativecontrol}
\end{equation}
Thus dissipativity does not sign the fermionic amplitudes.

There is also a scoped large-activity obstruction: a zero $w_0$ of $F$
inside $|w|<v$ gives an eigenvector of $S$ with eigenvalue $w_0/v$, hence
a kernel of $F(vS)$. At $|w_0|=v$, truncated unit-modulus eigenvectors
give approximate null vectors, so a bounded inverse still fails.
This limits the particular bounded-invertible generator formula; it
is not an exclusion of unbounded evolutions or enlarged systems.

\section{The fermionic loss instrument}
\label{sec:instrument}

\subsection{Finite modes and exact no-count evolution}
For $\mathfrak h_N=\operatorname{span}\{\ket0,\ldots,\ket N\}$ use the
restrictions $K_{v,N}$ and $G_N$, and the compressions $\loss_N,H_N$.
The scalar $F(v)$ remains complete even when $S_N$ is nilpotent. Put
\begin{equation}
 B_N=\sqrt{\loss_N},\quad L_{\mu,N}=\sum_j(B_N)_{\mu j}a_j,
 \quad H_{F,N}=d\Gamma(H_N),\quad \mathsf D_N=d\Gamma(G_N).
 \label{eq:lossrow}
\end{equation}
Normal ordering sets vacuum energy zero. The CAR give
\begin{align}
 \sum_\mu L_{\mu,N}^\dagger L_{\mu,N}&=d\Gamma(\loss_N),\notag\\
 \mathcal J_N\rho&=\sum_\mu L_{\mu,N}\rho L_{\mu,N}^\dagger
 =\sum_{j,k}(\loss_N)_{kj}a_j\rho a_k^\dagger,\notag\\
 \mathcal L_N\rho&=\mathsf D_N\rho+\rho\mathsf D_N^\dagger+
 \mathcal J_N\rho. \label{eq:lindblad}
\end{align}
The transposed coefficient $R_{kj}$ in the jump map is intentional.
The dissipator convention is $L\rho L^\dagger-\{L^\dagger L,\rho\}/2$,
as in the quantum-jump formulation of
Castin--Dalibard--M\o lmer \cite{CastinDalibardMolmer08054002v1}.
The finite-mode identities here are proved directly; no infinite-domain
result is imported from the trajectory literature.

Exterior differentiation proves the no-count operator
\begin{equation}
 M_{0,N}(b,a)=\Gamma(V_N(b,a)),\qquad
 \mathcal V_N(b,a)\rho=M_{0,N}(b,a)\rho M_{0,N}(b,a)^\dagger.
 \label{eq:nocount}
\end{equation}
For ordered count times $a<t_1<\cdots<t_q<b$, the instrument density is
\begin{equation}
 \mathcal I_{N;q}(\mathbf t;b,a)=
 \mathcal V_N(b,t_q)\mathcal J_N(t_q)\cdots
 \mathcal V_N(t_2,t_1)\mathcal J_N(t_1)\mathcal V_N(t_1,a).
 \label{eq:instrument}
\end{equation}
For $q=1$ the intervening factors have their evident meaning; $q=0$
is \eqref{eq:nocount}. Detector-channel labels are summed incoherently,
while count times are retained. Integrating these CP densities over
measurable record sets gives the count instrument. The series stops
at $q=r$ for inputs supported on at most $r$ particles. Its sum is
trace preserving by \eqref{eq:lindblad}, and splitting trajectories
at intermediate times proves composition. Thus the finite-mode
channel is CPTP and CP-divisible, generally time dependent.
These statements apply to arbitrary trace-class operators, including
coherences, rather than populations alone (\ev{O2,O3}).

\subsection{Hermitian collision realization and its topology}
At each fixed $N$, let $m=N+1$ and use $m$ fresh vacuum environment
modes in every time bin of length $h$. In environment-first order,
with environment parity $\Pi_E$, define
\begin{align}
 W_N&=\sum_\mu\left(b_\mu^\dagger\Pi_E\otimes L_{\mu,N}
             +\Pi_Eb_\mu\otimes L_{\mu,N}^\dagger\right),\notag\\
 U_{N,h}(v)&=\exp[-i(hI_E\otimes H_{F,N}(v)+\sqrt h\,W_N(v))].
 \label{eq:collision}
\end{align}
This exchange is Hermitian and conserves total number and parity.
Writing $C_\sigma=\bra\sigma U_{N,h}\ket{\Omega_E}$ gives, uniformly
on the local interval at fixed $N$,
\begin{equation}
 C_0=I+h\mathsf D_N+O_N(h^{3/2}),\quad
 C_\mu=-i\sqrt h L_{\mu,N}+O_N(h^{3/2}),\quad
 C_\sigma=O_N(h)\quad(|\sigma|\ge2). \label{eq:blocks}
\end{equation}
The signs follow from
$W_N(\Omega_E\otimes f)=\sum_\mu\ket\mu\otimes L_{\mu,N}f$ and
$\bra{\Omega_E}W_N^2\ket{\Omega_E}=d\Gamma(\loss_N)$.
Measuring each bin's total occupation produces the finite collision
count instrument.

\begin{theorem}[Fixed-mode collision limit, \ev{O2,O3}]
For a partition $\pi$ of $[a,b]$ with mesh $\delta\le1$, compare
collision counts with the exact Lindblad instrument integrated over
the \emph{same bins}. With $T_0=b-a$,
\begin{equation}
 \sup_{\|\rho\|_1\le1}\sum_{\mathbf n}
 \|\mathcal I^\pi_{{\rm coll},\mathbf n}\rho-
   \mathcal I^\pi_{{\rm Lind},\mathbf n}\rho\|_1
 \le T_0[(C_c+C_L)\sqrt\delta+\omega_N(\delta)]\longrightarrow0.
 \label{eq:recorderror}
\end{equation}
Here $C_c,C_L$ are explicit finite constants at fixed $N$, given in
Appendix~\ref{sec:constants}, and
\[
 \omega_N(\delta)=\sup_{|t-t'|\le\delta}
 \{2\|\mathsf D_N(t)-\mathsf D_N(t')\|
       +\|\mathcal J_N(t)-\mathcal J_N(t')\|_{1\to1}\}\to0.
\]
The ordered product of vacuum blocks also obeys the phase-sensitive bound
\begin{equation}
 \left\|\prod_l C_{0,l}-\Gamma(V_N(b,a))\right\|
 \le T_0[(E+d^2e^d/2)\sqrt\delta+\omega_D(\delta)]\to0,
 \label{eq:phaseerror}
\end{equation}
where $d$ bounds $\|\mathsf D_N\|$ and $\omega_D$ is its continuity modulus.
\end{theorem}
\begin{proof}[Reason for the error and the choice of record topology]
Taylor expansion of the unitary gives a local collision-instrument
error $C_ch^{3/2}$ against its first-order CP blocks. The frozen
Lindblad count instrument differs from those blocks by $C_Lh^2$.
A count register and Duhamel's formula bound time freezing by
$h\omega_N(\delta)$. Telescope the actual trace-preserving record
maps, whose induced trace norms are one, to obtain
\eqref{eq:recorderror}. The auxiliary first-order blocks are not used
as telescoping propagators. Telescoping the actual contraction
vacuum blocks separately gives \eqref{eq:phaseerror}; convergence of
a channel alone would not fix a Kraus phase.
\end{proof}

Assigning grid counts to bin representatives then gives weak
convergence of event-time tests. For a bounded continuous test on the
closed ordered simplex, the further error is its modulus of continuity
at $\delta$; on the open simplex, first exhaust away from boundaries
and collisions. There is no total-variation convergence claim between
atomic grid timestamps and a continuous timestamp law.

The duration-$h$ Hamiltonian is
$I_E\otimes H_{F,N}+h^{-1/2}W_N$. Fresh bins and singular short-time
coupling are resources of the limit; the constants may grow with $N$.
This is not an exact finite-bin device or a simultaneous infinite-mode
microscopic limit.

\subsection{A scoped obstruction to a regular fixed reservoir}
A locally bounded total Hamiltonian with an initially independent
fixed vacuum has $U(v,0)=I+O(v)$. Leakage from that vacuum subspace
therefore has probability $O(v^2)$. The prescribed no-count transfer,
for a normalized system vector, has
\begin{equation}
 \|K_{v,N}\psi\|^2
 =1-v\ip\psi{\loss_N(0)\psi}+O(v^2). \label{eq:shorttime}
\end{equation}
For example, $\psi=\ket0$ has coefficient $2\lambda_1>0$.
The nonzero linear loss contradicts that regular fixed-vacuum model.
This excludes precisely the initial-product, locally bounded,
fixed-vacuum realization class (\ev{A1,O3}); it does not exclude
endpoint dilations, singular couplings, initially correlated reservoirs,
or every non-Markovian model. The collision stream above lies outside
that class.

\section{Reading the signed amplitude}
\label{sec:readout}
A transition probability loses the sign of \eqref{eq:minoramp}.
Instead use a neutral flag and equal-particle-number branches,
\begin{equation}
 \Psi_{\rm in}=\frac{\ket0\ket I+\ket1\ket J}{\sqrt2},
 \qquad |I|=|J|=r. \label{eq:flaginput}
\end{equation}
Control both $H_{F,N}$ and every $L_{\mu,N}$ by the flag projector
onto $\ket1$. The reference arm is unchanged. The no-count block is
$P_0^{\rm flag}\otimes I+P_1^{\rm flag}\otimes\Gamma(K_{v,N})$.
Let $A=A_v(I,J)$ and
\begin{equation}
 p_J=\|\Gamma(K_{v,N})\ket J\|^2
     =\det[(K_{v,N}^\dagger K_{v,N})_{J,J}],\qquad
 q_0=\frac{1+p_J}{2}. \label{eq:success}
\end{equation}
The unnormalized flag density after no count is
\begin{equation}
 \frac12\begin{pmatrix}1&A^*\\A&p_J\end{pmatrix}. \label{eq:flagmatrix}
\end{equation}
Consequently
\begin{align}
 P(0,\pm_X)&=\frac{1+p_J\pm2\operatorname{Re}A}{4},&
 \operatorname{Re}A&=q_0\E[X\mid0],\notag\\
 P(0,\pm_Y)&=\frac{1+p_J\pm2\operatorname{Im}A}{4},&
 \operatorname{Im}A&=q_0\E[Y\mid0], \label{eq:readout}
\end{align}
where $Y=\bigl(\begin{smallmatrix}0&-i\\i&0\end{smallmatrix}\bigr)$.
For the real theta convention the $X$ difference is the signed target.
Signal conditioning uses $p_J$, while flag conditioning uses $q_0$;
they are not interchangeable. The inverse bound gives
\begin{equation}
 p_J\ge(12/13)^{2r}. \label{eq:survival}
\end{equation}
There is no positive signal-success floor uniform in $r$, and no
efficient-estimation conclusion. At $r=0$, $A=p_J=q_0=1$; at $v=0$,
$A=\delta_{I,J}$ and $p_J=q_0=1$.

Both arms have the same number and parity, including odd $r$.
A calibrated preparation compensates any known wedge-order sign.
No count leaves the same environment vacuum on the two arms;
\eqref{eq:phaseerror} preserves the coherent phase in the collision
limit. These controlled couplings and ideal counting are extra resources,
not properties available from an arbitrary channel black box.
Unrestricted dense mode couplings are allowed; no spatial-locality
or hardware implementation theorem is claimed.

For the polynomial control \eqref{eq:negativecontrol}, the very same
experiment returns
\begin{equation}
 A_{\rm pol}=-v^2/F_*(v)^2,\qquad
 p_*=(1+v^2+v^4)/F_*(v)^4,\qquad
 P(0,+_X)-P(0,-_X)=A_{\rm pol}<0. \label{eq:controlreadout}
\end{equation}
The physical construction therefore faithfully retains a negative
amplitude instead of converting it to a positive norm.

\section{Exact cutoff compatibility and vacuum Fock space}
\label{sec:extension}
For $M\ge N$, $G_M$ preserves $\mathfrak h_N$ and restricts to $G_N$.
For an operator supported on $\Fock(\mathfrak h_N)$, a jump term
$a_j\rho a_k^\dagger$ vanishes if either index exceeds $N$; the retained
block of $\loss_M$ is $\loss_N$. Hence the no-count maps, the count
instrument with channel labels discarded, and the full channel
restrict \emph{exactly} across cutoffs, by their bounded equations.
Neither $\sqrt{\loss_N}=P_N\sqrt\loss P_N$ nor compatibility of the
individual labeled jumps or finite collision unitaries is asserted.
The Hamiltonian alone need not preserve a cutoff; the no-count
cancellation in $G$ does.

For a particle cap $Q_r$ onto at most $r$ particles, let
$\mathfrak T_r=Q_r\mathfrak TQ_r$ and define the annihilation row
\begin{equation}
 \mathfrak a f=\sum_{j\ge0}e_j\otimes a_jf,
 \qquad \|\mathfrak a f\|^2=\ip f{\mathcal N f}\le r\|f\|^2.
\end{equation}
Then
\begin{equation}
 \mathcal J_{\loss}\rho=\Tr_{\mathfrak h}
 [ (\sqrt\loss\otimes I)\mathfrak a\rho\mathfrak a^\dagger
   (\sqrt\loss\otimes I)] \label{eq:infinitejump}
\end{equation}
is CP, lowers the cap, and has dual
$\mathfrak a^\dagger(\loss\otimes X)\mathfrak a$. Therefore
\begin{equation}
 \|\mathcal J_{\loss}-\mathcal J_{\loss'}\|_{1\to1}
 \le r\|\loss-\loss'\|,
 \qquad \|\mathcal L(v)\|_{\mathfrak T_r\to\mathfrak T_r}
 \le2r\|G\|+r\|\loss\|\le8rA_*. \label{eq:capbound}
\end{equation}
This supplies a norm-continuous bounded generator on each cap, without
requiring a spectral gap or a uniform square-root Lipschitz estimate.

\begin{theorem}[Compatible vacuum-Fock extension, \ev{O2,O3}]
The finite-mode count instruments extend uniquely to a strongly
continuous two-parameter CPTP evolution on the trace class of
$\Fock(\ell^2(\N_0))$, with the full no-count contraction
$\Gamma(V(b,a))$. If $\Pi_N$ projects onto $\Fock(\mathfrak h_N)$,
then for every trace-class $\rho$,
\begin{equation}
 \|\mathcal T(b,a)\rho-
 \mathcal T_N(b,a)(\Pi_N\rho\Pi_N)\|_1
 \le\|\rho-\Pi_N\rho\Pi_N\|_1\longrightarrow0, \label{eq:cutoff}
\end{equation}
uniformly over the allowed endpoints. The same estimate holds in
record-valued trace norm, and with particle-cap compressions.
\end{theorem}
\begin{proof}[Construction]
The projections $\Pi_N$ increase strongly to identity, and two-sided
compression converges in trace norm. Exact compatibility and
contractivity extend the finite-mode maps from their dense union.
Complete positivity passes at each finite amplification; trace
preservation, composition and statewise time continuity pass by density.
For instruments the target is the direct sum of trace-class-valued
$L^1$ spaces on ordered time simplexes, including the zero-count atom;
the same contraction argument applies. Applying contraction to
$\rho-\Pi_N\rho\Pi_N$ proves \eqref{eq:cutoff}.
\end{proof}

The dense invariant union $\bigcup_r\mathfrak T_r$ is a common
differentiability domain, as follows from the bounded ODE on each cap.
No maximal-domain or graph-core statement is made. Generator bounds
grow with $r$; the full-Fock derivative and inverse of $\Gamma(V)$
need not be bounded. The cutoff estimates are statewise, with no rate
without an input-tail bound and no all-state operator-norm convergence.
Physical density matrices respect fermion parity; the mathematical
extension to odd operators does not authorize parity-forbidden
preparations. The microscopic limit was taken at fixed $N$ in
Section~\ref{sec:instrument}; this separate trace-class extension
neither interchanges those limits nor proves an infinite-density
reservoir realization.

\section{Why the tested local sign mechanism stops}
\label{sec:cone}
\noindent\textbf{Evidence status.} The following obstruction is the
written B02 argument \ev{S1}, with proof-author checking \ev{S2}.
Its separate E04 verification remains pending. Reproducing it here does
not upgrade that status or close the broader MC08 investigation.

Let $C_{r,N}$ be the nonnegative real span of the increasing occupation
basis in $\wedge^r\mathfrak h_N$. A proposed path argument would require
\begin{equation}
 \wedge^rV_N(b,a)C_{r,N}\subseteq C_{r,N}
 \quad\hbox{for every intermediate }a\le b. \label{eq:conecondition}
\end{equation}
The infinitesimal generator must then have nonnegative off-diagonal
entries. This is the elementary differential-system obstruction behind
classical total-positivity results; see the compound-matrix discussion
and Theorems 6--7 in Margaliot--Sontag
\cite{MargaliotSontag180209590v2}. The argument below is finite-state
CAR algebra and does not import an infinite-dimensional theorem.

\begin{proposition}[Fixed-cone obstruction]
Within the analytic Toeplitz family \eqref{eq:G}, with $\lambda_1>0$,
condition \eqref{eq:conecondition} for $r=3$ and all cutoffs forces
$F(w)=c_0e^{\lambda_1w}$. The unchanged theta source fails this condition.
\end{proposition}
\begin{proof}
For $n\ge2$, $N\ge n+2$, set $\mathsf D_3=d\Gamma_3(G)$. Exactly one
hop contributes to each of
\begin{align}
 \bra{0,n+1,n+2}\mathsf D_3(v)\ket{n,n+1,n+2}&=+\alpha_n(v),\notag\\
 \bra{0,1,n+2}\mathsf D_3(v)\ket{1,n,n+2}&=-\alpha_n(v).
 \label{eq:oppositesigns}
\end{align}
The second hop crosses one occupied site. Since
$\wedge^3V_N(a+h,a)=I+h\mathsf D_3(a)+o(h)$, nonnegative off-diagonal
entries force $\alpha_n(a)=0$ for all $n\ge2$ at positive $a$.
Thus $\lambda_n=0$ for $n\ge2$, and analyticity gives the exponential.
Conversely that exponential has only adjacent allowed hops, all with
positive CAR sign, and preserves every finite ordered cone.

To rule out this exponential for theta, the full source obeys
\begin{equation}
 0<\Phi(u)<3000\exp(9u-\tfrac32e^{4u}),\qquad u\ge0.
 \label{eq:supergauss}
\end{equation}
Indeed split $\pi m^2e^{4u}\ge(\pi/2)e^{4u}+(\pi/2)m^2$, discard
only the negative subterm for an upper bound, and use
$\sum_{m\ge1}m^4e^{-(\pi/2)m^2}<\sum m^4 2^{-m}=150$ and $\pi^2<10$.
For every $A>0$ there is therefore a finite $C_A$ with
$\Phi(u)\le C_Ae^{-Au^2}$. Consequently, for $w\ge0$,
\begin{equation}
 F(w)\le\frac{C_A\sqrt\pi}{2\sqrt A}e^{w/(4A)},
 \qquad \lim_{w\to\infty}\frac{\log F(w)}w=0. \label{eq:nongrowth}
\end{equation}
The lower limit follows from $F(w)\ge c_0>0$; the upper limit follows
by taking arbitrary $A$. This contradicts the positive limiting slope
$\lambda_1>1/1000$ of $c_0e^{\lambda_1w}$.
\end{proof}

The failure has a finite intermediate witness. Let
$n_* =\min\{n\ge2:\lambda_n\ne0\}$, whose finiteness follows from the
preceding argument; no numerical value is asserted. Put
\begin{equation}
 B_*=[96(99/100)^3]^{-1},\qquad
 \mathcal M_3=3B_*+36A_*^2.
\end{equation}
The coefficient bounds imply $\|G'\|\le B_*$ and, on the three-particle
sector,
\[
 \|\wedge^3V(a+h,a)-I-h\mathsf D_3(a)\|
 \le\tfrac12\mathcal M_3h^2.
\]
For any $0<a<1/25$, choose the negative entry in
\eqref{eq:oppositesigns} and
\begin{equation}
 0<h\le\min\{(1/25-a)/2,\ |\alpha_{n_*}(a)|/\mathcal M_3\}.
 \label{eq:witness}
\end{equation}
Its intermediate amplitude is at most
$-h|\alpha_{n_*}(a)|/2<0$, already for $N\ge n_*+2$.
The constants do not depend on that cutoff.

The quantifiers matter: a fixed ordered cone, the unchanged generator,
all occupation inputs and every intermediate interval were required.
The argument does not disprove endpoint positivity from $a=0$.
For example $F_+(w)=1+w$ has nonnegative endpoint minors, while
\begin{equation}
 [K_bK_a^{-1}]_{0,2}
 =-a(b-a)\frac{1+a}{1+b}<0\qquad(0<a<b<1). \label{eq:bidiagonal}
\end{equation}
At $v=0$ the longer-hop coefficients vanish even for theta. Different
auxiliary flows, endpoint factorizations and cancellation among signed
paths remain outside the obstruction.

\section{Activity response cancels from the target ratio}
\label{sec:activity}
For $J_k=(k,k+1,k+2)$ and $I=(0,1,2)$,
\begin{equation}
 A_k(v)=\frac{v^{3k}}{F(v)^3}D_{3,k},\qquad
 \frac{A_{k-1}A_{k+1}}{A_k^2}=\eta_k,
 \quad k\ge1,\quad 0<v\le1/25. \label{eq:invariant}
\end{equation}
The denominators are positive by the inherited order-three result.
The physical ratio is exactly independent of activity, not just
approximately stationary.

To see how conditioning enters, set
$\psi_k=\wedge^3K_v\ket{J_k}$, $p_k=\|\psi_k\|^2$ and
$\ell_k=\ip{\psi_k}{d\Gamma_3(\loss)\psi_k}\ge0$. Then
$p_k'=-\ell_k$, $q_k=(1+p_k)/2$, and the signed readout is
$A_k=q_kx_k$ with $x_k=\E[X\mid0,J_k]$. For $v>0$,
\begin{align}
 \partial_v\log A_k&=3k/v-3\gamma(v),\notag\\
 \partial_v\log q_k&=-\ell_k/(1+p_k),\notag\\
 \partial_v\log x_k&=3k/v-3\gamma(v)+\ell_k/(1+p_k).
 \label{eq:responsecancel}
\end{align}
Restoring the success factor cancels the loss term exactly, and the
neighboring second difference cancels the remaining activity terms:
\begin{equation}
 \partial_v\log\left(
 \frac{q_{k-1}q_{k+1}}{q_k^2}\frac{x_{k-1}x_{k+1}}{x_k^2}
 \right)=\partial_v\log\eta_k=0. \label{eq:ratioresponse}
\end{equation}
The initial state $K_0=I$ does not give a useful ratio anchor:
$A_k(0)=0$ for $k\ge1$, and its limiting ratio is precisely the unknown
coefficient ratio. Increasing the admissible activity interval cannot
change it. A useful information or response inequality would therefore
have to couple the three \emph{signed endpoint values} after complete
reconstruction. These identities do not exclude every such inequality;
they identify why separate survival or activity-response bounds have
not supplied one (\ev{S1,S2}).

\section{The natural theta energies and what the orbitals measure}
\label{sec:energy}

\subsection{Circle energies, the primitive and its positive filter}
The periodic circle Hamiltonian $L=-\pi\partial_x^2$ has eigenvectors
$e^{imx}/\sqrt{2\pi}$ and eigenvalues $\pi m^2$, $m\in\mathbb Z$.
Thus the \href{https://dlmf.nist.gov/20.2.E3}{defining theta series} \cite[20.2.3]{NISTDLMFPositiveThetaPrimitives}
gives
\begin{equation}
 Q(\beta)=\sum_{m\ge1}e^{-\beta E_m}
 =\tfrac12(\Tr e^{-\beta L}-1),\qquad E_m=\pi m^2,
 \qquad p(u)=e^uQ(e^{4u}). \label{eq:energies}
\end{equation}
The subtraction removes the zero mode and $1/2$ removes the twofold
$\pm m$ degeneracy. Equivalently $Q$ is the trace on the positive-momentum
subspace. These energy units match the canonical exponent; no measured
energy scale is being inferred.

Since $\beta'=4\beta$,
$p''=e^u(Q+24\beta Q'+16\beta^2Q'')$, and hence
\begin{align}
 \Phi(u)&=e^u\sum_{m\ge1}P_0(\beta E_m)e^{-\beta E_m}
 =\frac{e^u}{2}\Tr[P_0(\beta L)e^{-\beta L}],\notag\\
 P_0(x)&=x(2x-3). \label{eq:filter}
\end{align}
On $u\ge0$, every nonzero $x=\beta E_m$ satisfies $x\ge\pi>3/2$.
The filter is positive and the zero mode contributes zero. The
polynomial times heat damping is trace class. The normalized
positive-mode distribution is
\begin{equation}
 \rho_u(m)=\frac{P_0(\beta E_m)e^{-\beta E_m}}
                 {\sum_{n\ge1}P_0(\beta E_n)e^{-\beta E_n}}.
 \label{eq:thermalstate}
\end{equation}
It is a filtered thermal state, with a temperature-dependent insertion,
not the original Gibbs state. Filter positivity is asserted here only
for $\beta\ge1$. For the ordinary positive-mode Gibbs expectation,
\begin{equation}
 \frac{\Phi(u)}{p(u)}=2\beta^2\langle E^2\rangle_\beta
                          -3\beta\langle E\rangle_\beta. \label{eq:energymean}
\end{equation}
Full theta inversion also gives $p(-u)=p(u)+\sinh u$; applying
$\partial_u^2-1$ recovers the complete source's evenness. These
identities retain every spectral mode (\ev{T1,T2}).

\subsection{Response order is not energy level}
Define the exact polynomials
\begin{equation}
 P_{n+1}(x)=4xP_n'(x)+(1-4x)P_n(x). \label{eq:recurrence}
\end{equation}
Termwise differentiation of \eqref{eq:filter} yields
\begin{equation}
 \frac{\Phi^{(n)}(u)}{\Phi(u)}
 =\E_{\rho_u}\frac{P_n(X)}{P_0(X)},\qquad X=\beta E_m.
 \label{eq:responseenergy}
\end{equation}
Thus the original derivative orbital is exactly
\begin{equation}
 f_i(u)=\sqrt{\Phi(u)}\,
       \E_{\rho_u}\frac{P_{2i}(X)}{P_0(X)}. \label{eq:orbitalenergy}
\end{equation}
The label $i$ selects an energy-response observable; $m$ selects the
energy eigenmode being averaged. Every $f_i$ contains the complete
spectrum. Likewise $c_j$ integrates all levels over $u$ with its
moment weight, and is not a single Boltzmann factor for $E_j$.

At order three the required polynomials are
\begin{align}
 P_2(x)&=32x^4-224x^3+330x^2-75x,\notag\\
 P_4(x)&=512x^6-8448x^5+41408x^4-68096x^3
                  +30930x^2-1875x. \label{eq:polys}
\end{align}
Put $r_2=P_2/P_0$, $r_4=P_4/P_0$. Multilinearity gives
\begin{equation}
 \det\begin{bmatrix}1&1&1\\\srcR(u_1)&\srcR(u_2)&\srcR(u_3)\\
 \srcS(u_1)&\srcS(u_2)&\srcS(u_3)\end{bmatrix}
 =\E_{\rho_{u_1}\otimes\rho_{u_2}\otimes\rho_{u_3}}
 \det\begin{bmatrix}1&1&1\\r_2(X_1)&r_2(X_2)&r_2(X_3)\\
 r_4(X_1)&r_4(X_2)&r_4(X_3)\end{bmatrix}. \label{eq:threecopy}
\end{equation}
This is a faithful spectral representation of the derivative determinant,
not a positivity theorem. Ordering $u_1<u_2<u_3$ does not order
$X_j=\pi e^{4u_j}m_j^2$ for every independently sampled occupation tuple.

\subsection{A genuine energy-mode Fock Hamiltonian}
On the positive-mode space one may define independent fermions
$\mathfrak b_m$ with
\begin{equation}
 \mathbb H_\theta=\sum_{m\ge1}E_m\mathfrak b_m^\dagger\mathfrak b_m,
 \qquad E_J=\sum_{m\in J}E_m. \label{eq:energyfock}
\end{equation}
It is the nonnegative occupation-energy operator on its natural domain.
For $\beta>0$ its grand partition function is
\begin{equation}
 \Xi(\beta,z)=\prod_{m\ge1}(1+ze^{-\beta E_m}),\qquad
 [z]\Xi(\beta,z)=Q(\beta). \label{eq:grand}
\end{equation}
The product follows by summing zero or one occupation of each mode;
$\sum_m e^{-\beta E_m}<\infty$ gives convergence locally uniformly in
$z$. The primitive uses the \emph{one-particle coefficient} of this
product, not the entire grand partition function. One-particle heat
traces themselves do not distinguish particle statistics.

The energy creators $\mathfrak b_m^\dagger$ are not identified with
coefficient-site creators $a_i^\dagger$. Moreover, the MC08 Hamiltonian
$H(v)=i(G-G^\dagger)/2$ is engineered from the transfer and generally
mixes sites; it is not the diagonal $L$. Attaching the energies
$\pi i^2$ to coefficient sites without a proved relation would change
the model. The exact link already available is the observable identity
\eqref{eq:orbitalenergy} (\ev{D2}); an operator correspondence with
additional sign consequences remains to be found.

\section{The signed thermal response that remains uncontrolled}
\label{sec:residual}
Write $\ell_j=\partial_u^j\log\Phi$. The temperature dependence of the
filter is essential. For one summand let
\begin{equation}
 a(u,m)=u+\log P_0(x)-x,\quad x=\beta E_m,\quad
 q(x)=\partial_u a=9-4x+\frac{12}{2x-3},
\end{equation}
and set $b=4xq'$, $c=4xb'$, $d=4xc'$. In particular
$b=-16x-96x/(2x-3)^2<0$ for $x\ge\pi$.
For an energy observable $f(X)$,
\begin{equation}
 \partial_u\E_{\rho_u}f=\E_{\rho_u}(4Xf')+
                  \Cov_{\rho_u}(f,q). \label{eq:thermderivative}
\end{equation}
Repeated differentiation gives the complete identities
\begin{align}
 \ell_1&=\E q,&
 \ell_2&=\E b+\Var(q),\notag\\
 \ell_3&=\E c+3\Cov(q,b)+\kappa_3(q),\notag\\
 \ell_4&=\E d+4\Cov(q,c)+3\Var(b)
                 +6\kappa(q,q,b)+\kappa_4(q). \label{eq:cumulants}
\end{align}
Here $\kappa_3(q)$ is the third centered moment,
$\kappa(q,q,b)=\E[(q-\E q)^2(b-\E b)]$, and
$\kappa_4(q)=\E[(q-\E q)^4]-3\Var(q)^2$.
These are thermal cumulants, distinct from the comparator
$\kappa_3(s)$ in \eqref{eq:curvtarget}. The positive variance in
$\ell_2$ opposes the negative individual-mode term, and higher mixed
cumulants have no sign supplied by positivity of the state.

For the explicit single-energy comparison
$\phi_*(u)=e^uP_0(\pi e^{4u})e^{-\pi e^{4u}}$, the energy distribution
is a point mass, so its thermal fluctuations vanish. Nevertheless, at
$x=4$ the exact response-polynomial calculation gives
\begin{equation}
 r_2'(4)=\frac{952}{25}>0,\qquad
 \frac12\frac{d^2r_4}{dr_2^2}(4)
 =-\frac{157375}{99127}<0. \label{eq:thermalwitness}
\end{equation}
The second number is the confluent three-node determinant quotient.
Thus a positive filtered state, zero thermal variance and individual-mode
log-concavity do not imply pointwise positivity of the derivative
observable. This comparison lacks the complete theta spectrum and
reflection; it is not a counterexample to the canonical integrated
comparison. The actual theta insertion already has a negative origin
region by the background analysis (\ev{T1,T2,BG}).

The residual can be displayed exactly. Let
\begin{equation}
 U=\srcS-\srcR^2
 =\ell_4+4\ell_1\ell_3+2\ell_2^2+4\ell_1^2\ell_2
 =\srcR''+2\ell_1\srcR'. \label{eq:U}
\end{equation}
The background supplies $\srcR'>0$ on the interior, so regard $U$ as a
function of $y=\srcR(u)$. If $\chi(\mathbf u)=U[y_1,y_2,y_3]$ is its
second divided difference, then
\begin{equation}
 h_3=1+\chi. \label{eq:chi}
\end{equation}
For the Gaussian source $e^{-au^2}$, $U=-8a\srcR-8a^2$ is affine,
so $\chi=0$ and $h_3=1$. This is a comparator normalization, not a
bound for theta.

For clarity, the positive carrier recalled from the background is
\begin{equation}
 w_s(\mathbf u)=\Bigl(\prod_{j=1}^3u_j^{2s}\Phi(u_j)\Bigr)
               V(\mathbf u^2)V(\srcR(\mathbf u)),\quad
 d\mu_s=Z(s)^{-1}w_s\dd\mathbf u,\quad \mathbf u\in\Omega_3.
 \label{eq:carrier}
\end{equation}
Here $\Omega_3=\{0<u_1<u_2<u_3\}$, $V$ is the Vandermonde, and
$M_3=Z\mathcal A$ with $\mathcal A=1+\E_s\chi>0$.
Let $T=2\sum_j\log u_j$ and $m_s=\E_sT$. The complete curvature is
\begin{equation}
 (\log M_3)''=
 \Var_sT+\frac{\Cov_s(\chi,(T-m_s)^2)}{\mathcal A}
 -\left(\frac{\Cov_s(\chi,T)}{\mathcal A}\right)^2. \label{eq:completecurv}
\end{equation}
Thus the missing comparison \eqref{eq:curvtarget} is exactly
\begin{equation}
 \mathcal A^2[\kappa_3(s)-\Var_sT]
 -\mathcal A\Cov_s(\chi,(T-m_s)^2)
 +\Cov_s(\chi,T)^2\ge0. \label{eq:unpaid}
\end{equation}
Equation \eqref{eq:unpaid} is a restatement of the target, not an acquired
estimate. It combines energy responses at each source point, a divided
difference across three points, and integration against the Mellin
carrier. A variance bound at one temperature does not control all three
operations.

For each fixed derivative order $n$, the complete series has an envelope
$|\Phi^{(n)}(u)|\le C_ne^{(4n+9)u}e^{-\pi e^{4u}}$ on $u\ge0$.
This justifies the fixed-order thermal differentiations and log-inserted
integrals. At node collisions use the smooth undivided product
$h_3w_s=(\prod u_j^{2s})V(\mathbf u^2)\det[\Phi^{(2i)}(u_j)]$;
near zero, $u^{2s}|\log u|^p$ is locally integrable for $s\ge0$.
These are absolute-integrability payments, not determinant-relative
estimates or a large-$s$ comparison theorem.

\section{What is established and what the next relation must add}
\label{sec:next}
The outcome is a faithful physical formulation of the original signed
problem. The loss model realizes its amplitudes, preserves their
normalization and phase, and has stated finite/infinite limits.
The energy description identifies what the derivative labels measure.
Neither representation supplies the missing signed comparison.

The exact Gaussian benchmark is
\begin{equation}
 D^G_{3,k}=\frac{2c_0^3t^{3k}}{k!(k+1)!(k+2)!},\qquad
 \frac{D^G_{3,k-1}D^G_{3,k+1}}{(D^G_{3,k})^2}=\frac{k}{k+3}.
\end{equation}
Thus log-concavity of $D_{3,k}/D^G_{3,k}$ is an equivalent rewriting
of \eqref{eq:target}. Likewise
\begin{equation}
 W_k(v)=\frac{k}{k+3}A_k(v)^2-A_{k-1}(v)A_{k+1}(v)\ge0
 \label{eq:jointtarget}
\end{equation}
is a useful joint expression, but merely naming $W_k$ does not prove it.

Three possibilities discussed in \ev{D1,D2} remain conceptual:
\begin{enumerate}
\item A source-derived pairing or exchange identity controlling
\eqref{eq:jointtarget}, with signed contributions combined before any
bound is applied.
\item An integrated relation among the complete energy responses
\eqref{eq:responseenergy} that bounds \eqref{eq:completecurv}, or only
its necessary triangular averages. The square spectrum and full theta
reflection are available information beyond generic thermal positivity.
\item An auxiliary Gaussian-to-theta deformation of the \emph{source},
with positivity along the path and a paid derivative comparison. This
would be a new flow, not variation of the activity $v$.
\end{enumerate}
No mechanism from this list is selected or proved by the manuscript.
The energy-response relation is the author's current discussion focus.
A proposed estimate should identify the signed quantity it controls,
its source-specific premise, its actual domain and its treatment of
boundaries and tails. Existing MC03 localization, MC04 Ward-response
nonclosure and MC05 higher-coefficient nonclosure, summarized in the
background, must remain part of that assessment.

\begin{table}[htbp]
\centering\small
\renewcommand{\arraystretch}{1.17}
\begin{tabularx}{\textwidth}{@{}>{\raggedright\arraybackslash}p{.24\textwidth}X@{}}
\toprule
Control or scope & What it actually establishes\\
\midrule
Gaussian generating function & Passes local cone preservation and saturates the benchmark; it is not the canonical theta source.\\
$1+w+2w^2$ & Passes dissipativity and the same readout, but has a negative endpoint minor. It fails canonical source reconstruction.\\
$1+w$ & Positive endpoint minors coexist with a negative intermediate transfer entry.\\
One-energy thermal source & Positive filtered state and zero fluctuations coexist with a negative derivative observable; full spectrum and reflection fail.\\
Weighted-square / Euler changes & Change canonical weights or reconstruction and can lose reflection/endpoints. Their uncomputed signs and local admissibility are not assigned here.\\
Reflected / shape controls & Can retain selected symmetry or decay while lacking the complete arithmetic expansion. Such shape information alone is not a signed comparison.\\
\bottomrule
\end{tabularx}
\caption{The controls distinguish the premises of a realization from the missing signed estimate.}
\end{table}

The stop concerns one every-intermediate-interval fixed-cone mechanism,
not the whole subinvestigation. Activity cancellation leaves an unpaid
joint endpoint relation; it does not exclude every quantum-information
method. The next session can begin from the exact spectral dictionary
and ask for that additional relation, without rebuilding the loss model
or replacing the original observable by a positive probability.

\appendix
\section{Explicit fixed-mode collision constants}
\label{sec:constants}
This appendix retains the quantitative content of \ev{O2}, equations
(8)--(16), without claiming constants uniform in the cutoff. For
$m=N+1$, $A=A_*$, choose
\[
 d=H_0=2mA,\quad j=4mA,\quad \ell=2\sqrt A,\quad
 w=4m\sqrt A,\quad c=w+H_0,
\]
bounding respectively the no-count generator, Fock Hamiltonian,
jump map, each loss operator, and exchange. Set
\begin{align}
 E&=wH_0+\tfrac12H_0^2+e^cc^3/6,\qquad q=2^m-m-1,\notag\\
 C_c&=E(2+d)+mE(2\ell+E)+q(w^2/2+E)^2,\notag\\
 C_L&=(2+d)d^2e^d/2+2dj+j^2e^j/2. \label{eq:explicitconstants}
\end{align}
Expansion of $e^{-i(\sqrt hW+hH)}$ through degree two has remainder
at most $Eh^{3/2}$ for $0<h\le1$. The $q$ multiple-occupation blocks
contribute only order $h^2$ in the instrument norm. The three terms in
$C_L$ bound the frozen vacuum error, one-count integral error and the
sum of higher-count terms. For the time-dependent modulus,
\[
 \omega_N(\delta)\le m[2\omega_G(\delta)+\omega_{\loss}(\delta)]
 \longrightarrow0
\]
by \eqref{eq:capbound}. These constants prove
\eqref{eq:recorderror} and \eqref{eq:phaseerror} on arbitrary
trace-class inputs, including flag coherences. They are conservative
analytical bounds, not performance estimates for a laboratory device.

\section{Evidence map and continuation boundary}
\label{sec:evidence}
The source root for the following identifiers is
\path{papers/investigations-5/07_positive_theta_primitives/01_jensen/01_mellin_curvature/}.
Paths are repository-relative to that root. The build record freezes
file hashes; hashes identify versions and do not validate proofs.

\begin{description}[style=nextline,leftmargin=0pt,labelwidth=0pt,labelsep=0pt]
\item[\ev{BG} Background manuscript.]
\path{manuscript/main.tex}, \emph{Signed Mellin Curvature and Theta
Coefficient Positivity}, 2 November 2026; source normalization,
real-tilt positivity baseline, factorial budget and eight-branch context.
Its production self-check is exposition checking, not a new theorem review.
\item[\ev{F1,F2} Fermionic representation and review.]
\path{02_fermionic_correlators/notes/01_THETA_ORBITAL_COMPARISON_RETURN_20261101.md}
and \path{02_fermionic_correlators/reviews/01_MC02_E01_SEPARATE_REVIEW_20261101.md}.
The transfer, continuum overlap and fixed-positive-metric obstruction
have separate cross-model internal LLM review with shared premises.
\item[\ev{O1} Local admissibility.]
\path{08_open_quantum_dynamics/notes/01_NO_JUMP_ADMISSIBILITY_AUDIT_20261102.md},
with the exact rational support record in that branch's numerics folder.
The full-source proof is analytical, not a numerical theta cutoff.
\item[\ev{A1} Architecture decision.]
\path{08_open_quantum_dynamics/notes/02_OPEN_SYSTEM_ARCHITECTURE_DECISION_20261102.md}.
This fixes the fresh-vacuum loss/count resources and states the scoped
regular-reservoir obstruction; its construction is completed in O2.
\item[\ev{O2,O3} Realization and separate assessment.]
\path{08_open_quantum_dynamics/notes/03_FERMIONIC_LINDBLAD_REALIZATION_20261102.md}
and \path{08_open_quantum_dynamics/reviews/02_OPEN_QUANTUM_DYNAMICS_ASSESSMENT_20261102.md}.
O3 accepts O1 and O2 at their stated scope without correction, in a
separate same-model internal review. It does not supply a signed estimate.
\item[\ev{S1,S2} Scoped sign-mechanism stop and self-check.]
\path{08_open_quantum_dynamics/notes/04_SIGNED_COMPARISON_DECISION_20261102.md}
and \path{08_open_quantum_dynamics/reviews/03_MC08_B02_STOP_SELF_CHECK_20261102.md}.
These have proof-author checking only. The sole live progress register
narrows earlier parking language to the tested mechanism; the broader
folder remains open. E04 is Ready and unstarted at this manuscript freeze.
\item[\ev{T1,T2} Thermal identities and review.]
\path{01_circle_thermal_trace/notes/01_THERMAL_SOURCE_IDENTITY_RETURN_20261101.md}
and \path{01_circle_thermal_trace/reviews/01_MC01_E01_SEPARATE_REVIEW_20261101.md}.
The response polynomials, mixed cumulants, single-energy witness and
residual are accepted without correction in a separate cross-model
internal review. The full Mellin comparison remains unpaid.
\item[\ev{D1,D2} Current discussion.]
\path{08_open_quantum_dynamics/notes/05_MATHEMATICAL_RELATION_BRAINSTORM_20261102.md}
and \path{08_open_quantum_dynamics/notes/06_THETA_ENERGY_LEVELS_AND_FERMIONIC_LABELS_20261102.md}.
These organize candidate relations and distinguish the labels and
Hamiltonians; they are not new signed bounds.
\end{description}

The source/bibliography checks for this manuscript are selected-passage
checks, not a literature novelty audit. The accompanying production
self-check concerns transcription, scope, references and rendering; it
does not execute E04 or constitute a separate mathematical assessment.
The sole live register and the next-session note remain under
\path{08_open_quantum_dynamics/research-progress/}. The default next task
is \textbf{MC08-B03, conceptual discussion of a source-dependent
energy-response relation}; the separate E04 proof check remains an
explicit alternative. No numerical campaign is authorized by the
manuscript itself.

\section*{Acknowledgment and preparation record}
This manuscript was drafted for Edward Baker with substantial assistance
from OpenAI's \texttt{gpt-6-astra} model, recorded reasoning effort
\texttt{xhigh}, on 2 November 2026. It synthesizes research notes prepared
with the model configurations recorded in those notes. The present
preparation is an expository integration and author-side consistency
check; it does not replace independent mathematical review. No new
manuscript snapshot directory or scientific computation campaign was
created.

\bibliographystyle{plainurl}
\bibliography{manuscript_references}
\end{document}
